% GENERATED FILE. Edit canonical lessons, then run npm run generate:books.
% canonical-source-hash: fnv1a64:f8a95ca0130a0c95
% canonical-lessons: ZH-C187-dapo, ZH-C187-shuai, ZH-C187-houhui, ZH-C187-dangshi, ZH-C187-jieguo, ZH-R187-review-182-187, ZH-R187-review-final-187

\chapter{What a Day}
\label{ch:zh-what-a-day}
\hlchaptermodality{187}

\begin{chapteropening}
\textbf{By the end of this chapter:} \emph{I can tell what went wrong: something smashed, a fall, a regret, how things were then and how it ended.}

Complete the last lesson of chapter 187: I can tell what went wrong: something smashed, a fall, a regret, how things were then and how it ended.
\end{chapteropening}

\section[dǎpò]{dǎpò — to break; to smash}

\label{lesson:ZH-C187-dapo}



\begin{quote}
Before the new one: say the Mandarin for a bottle, then the Mandarin for a drawer.
\end{quote}

\subsection*{You'll want to know: dǎpò}
\emph{dǎpò} — \textquotedblleft{}to break; to smash\textquotedblright{}.

\begin{quote}
\emph{zuótiān wǒ dǎpò le yí ge bēizi} — Yesterday I broke a glass.
\end{quote}

Its characters come later; for now it is a word for the ear.

\begin{cousinweb}
\emph{dǎpò} is built from \zh{打} 'to strike' and \emph{pò} 'to break', both inherited from Middle Chinese, \emph{taengX} and \emph{phaH} (Guangyun, in Baxter's transcription), though dǎ is not the sound \emph{taengX} would regularly give.
\end{cousinweb}

\begin{grammarlens}[title={What you've built}]
One more word for talking about the past.
\end{grammarlens}

\subsection*{Your turn}
\begin{itemize}
\raggedright
  \item \emph{Say it:} \emph{dǎpò}
  \item \emph{Say it:} \emph{dǎpò}, once more
  \item \emph{Say it:} \emph{zuótiān wǒ dǎpò le yí ge bēizi} — Yesterday I broke a glass.
  \item \emph{Recall:} say the Mandarin for a bottle, then the Mandarin for a drawer, then say \emph{dǎpò} again
\end{itemize}

\begin{tcolorbox}[breakable,colback=teal!4,colframe=teal!35!black,title={Before you move on}]
What does dǎpò mean? (\textquotedblleft{}To break; to smash\textquotedblright{}.) Say it once more.
\end{tcolorbox}
\section[shuāi]{shuāi — to have a fall; to drop (and break)}

\label{lesson:ZH-C187-shuai}



\begin{quote}
Before the new one: say the Mandarin for a drawer, then the Mandarin for to smash.
\end{quote}

\subsection*{You'll want to know: shuāi}
\emph{shuāi} — \textquotedblleft{}to have a fall; to drop (and break)\textquotedblright{}.

\begin{quote}
\emph{tā shuāi le shǒu hěn téng} — He had a fall, and his hand really hurts.
\end{quote}

Its characters come later; for now it is a word for the ear.

\begin{cousinweb}
\emph{shuāi} No dictionary we use records an older form, so its history is not settled.
\end{cousinweb}

\begin{grammarlens}[title={What you've built}]
One more word for talking about the past.
\end{grammarlens}

\subsection*{Your turn}
\begin{itemize}
\raggedright
  \item \emph{Say it:} \emph{shuāi}
  \item \emph{Say it:} \emph{shuāi}, once more
  \item \emph{Say it:} \emph{tā shuāi le shǒu hěn téng} — He had a fall, and his hand really hurts.
  \item \emph{Recall:} say the Mandarin for a drawer, then the Mandarin for to smash, then say \emph{shuāi} again
\end{itemize}

\begin{tcolorbox}[breakable,colback=teal!4,colframe=teal!35!black,title={Before you move on}]
What does shuāi mean? (\textquotedblleft{}To have a fall; to drop (and break)\textquotedblright{}.) Say it once more.
\end{tcolorbox}
\section[hòuhuǐ]{hòuhuǐ — to regret}

\label{lesson:ZH-C187-houhui}



\begin{quote}
Before the new one: say the Mandarin for to smash, then the Mandarin for to have a fall.
\end{quote}

\subsection*{You'll want to know: hòuhuǐ}
\emph{hòuhuǐ} — \textquotedblleft{}to regret\textquotedblright{}.

\begin{quote}
\emph{wǒ hěn hòuhuǐ zuótiān méi qù wǎnhuì} — I really regret not going to the party yesterday.
\end{quote}

Its characters come later; for now it is a word for the ear.

\begin{cousinweb}
\emph{hòuhuǐ} is built from \zh{后} 'after' and \emph{huǐ} 'to repent', both inherited from Middle Chinese, \emph{huwX} and \emph{xwojH} (Guangyun, in Baxter's transcription), which glosses the second 'to repent and change'; today's third tone goes with its other reading, \emph{xwojX}. To repent afterwards.
\end{cousinweb}

\begin{grammarlens}[title={What you've built}]
One more word for talking about the past.
\end{grammarlens}

\subsection*{Your turn}
\begin{itemize}
\raggedright
  \item \emph{Say it:} \emph{hòuhuǐ}
  \item \emph{Say it:} \emph{hòuhuǐ}, once more
  \item \emph{Say it:} \emph{wǒ hěn hòuhuǐ zuótiān méi qù wǎnhuì} — I really regret not going to the party yesterday.
  \item \emph{Recall:} say the Mandarin for to smash, then the Mandarin for to have a fall, then say \emph{hòuhuǐ} again
\end{itemize}

\begin{tcolorbox}[breakable,colback=teal!4,colframe=teal!35!black,title={Before you move on}]
What does hòuhuǐ mean? (\textquotedblleft{}To regret\textquotedblright{}.) Say it once more.
\end{tcolorbox}
\section[dāngshí]{dāngshí — at that time; then}

\label{lesson:ZH-C187-dangshi}



\begin{quote}
Before the new one: say the Mandarin for to have a fall, then the Mandarin for to regret.
\end{quote}

\subsection*{You'll want to know: dāngshí}
\emph{dāngshí} — \textquotedblleft{}at that time; then\textquotedblright{}.

\begin{quote}
\emph{dāngshí wǒ bú zài jiā} — I wasn't at home at the time.
\end{quote}

Its characters come later; for now it is a word for the ear.

\begin{cousinweb}
\emph{dāngshí} is built from \emph{dāng} and \zh{时} 'time', both inherited from Middle Chinese, \emph{tang} and \emph{dzyi} (Guangyun, in Baxter's transcription). The dictionary glosses the first 'to match' and 'to meet with': the time it met with.
\end{cousinweb}

\begin{grammarlens}[title={What you've built}]
One more word for talking about the past.
\end{grammarlens}

\subsection*{Your turn}
\begin{itemize}
\raggedright
  \item \emph{Say it:} \emph{dāngshí}
  \item \emph{Say it:} \emph{dāngshí}, once more
  \item \emph{Say it:} \emph{dāngshí wǒ bú zài jiā} — I wasn't at home at the time.
  \item \emph{Recall:} say the Mandarin for to have a fall, then the Mandarin for to regret, then say \emph{dāngshí} again
\end{itemize}

\begin{tcolorbox}[breakable,colback=teal!4,colframe=teal!35!black,title={Before you move on}]
What does dāngshí mean? (\textquotedblleft{}At that time; then\textquotedblright{}.) Say it once more.
\end{tcolorbox}
\section[jiéguǒ]{jiéguǒ — in the end; as a result; a result}

\label{lesson:ZH-C187-jieguo}



\begin{quote}
Before the new one: say the Mandarin for to regret, then the Mandarin for at that time.
\end{quote}

\subsection*{You'll want to know: jiéguǒ}
\emph{jiéguǒ} — \textquotedblleft{}in the end; as a result; a result\textquotedblright{}.

\begin{quote}
\emph{wǒmen děng le yí ge xiǎoshí jiéguǒ tā méi lái} — We waited an hour, and in the end he didn't come.
\end{quote}

Its characters come later; for now it is a word for the ear.

\begin{cousinweb}
\emph{jiéguǒ} is built from \emph{jié} 'to tie' and \emph{guǒ} 'fruit', both inherited from Middle Chinese, \emph{ket} and \emph{kwaX} (Guangyun, in Baxter's transcription): to set fruit, and so an outcome. The step to 'in the end' is not given there.
\end{cousinweb}

\begin{grammarlens}[title={What you've built}]
One more word for talking about the past.
\end{grammarlens}

\subsection*{Your turn}
\begin{itemize}
\raggedright
  \item \emph{Say it:} \emph{jiéguǒ}
  \item \emph{Say it:} \emph{jiéguǒ}, once more
  \item \emph{Say it:} \emph{wǒmen děng le yí ge xiǎoshí jiéguǒ tā méi lái} — We waited an hour, and in the end he didn't come.
  \item \emph{Recall:} say the Mandarin for to regret, then the Mandarin for at that time, then say \emph{jiéguǒ} again
\end{itemize}

\begin{tcolorbox}[breakable,colback=teal!4,colframe=teal!35!black,title={Before you move on}]
What does jiéguǒ mean? (\textquotedblleft{}In the end; as a result; a result\textquotedblright{}.) Say it once more.
\end{tcolorbox}
\section[(dialogue)]{Looking back: from Eating Out to What a Day}

\label{lesson:ZH-R187-review-182-187}



\begin{quote}
Say the Mandarin for a restaurant and a cold starter.
\end{quote}

\subsection*{Your turn}
\begin{itemize}
\raggedright
  \item \emph{Say it:} a restaurant, a cold starter, a snack, home cooking, a boxed meal
  \item \emph{Say it:} to go out of here, to come inside, to finish work, to lock up, to lie down
  \item \emph{Say it:} to go on holiday, to decide, to look round a place, one's home town, to relax
  \item \emph{Say it:} for ages, before (a time), after (a time), morning and evening, more often than not
  \item \emph{Say it:} scissors, a rubbish bin, a plastic bag, a bottle, a drawer
  \item \emph{Say it:} to smash, to have a fall, to regret, at that time, in the end
  \item \emph{Read it:} \zh{冷} — each one aloud, then a word that holds it
\end{itemize}

\begin{tcolorbox}[breakable,colback=teal!4,colframe=teal!35!black,title={Before you move on}]
A restaurant? (\emph{fànguǎn}.) A cold starter? (\textbf{\zh{冷菜}}, \emph{lěngcài}.) A snack? (\textbf{\zh{小吃}}, \emph{xiǎochī}.) Home cooking? (\textbf{\zh{家常菜}}, \emph{jiāchángcài}.) A boxed meal? (\emph{héfàn}.) To go out of here? (\textbf{\zh{出去}}, \emph{chūqù}.) To come inside? (\textbf{\zh{进来}}, \emph{jìnlái}.) To finish work? (\emph{xiàbān}.) To lock up? (\emph{suǒmén}.) To lie down? (\emph{tǎng}.) To go on holiday? (\emph{dùjià}.) To decide? (\emph{juédìng}.) To look round a place? (\emph{cānguān}.) One's home town? (\emph{lǎojiā}.) To relax? (\emph{fàngsōng}.) For ages? (\emph{hǎojiǔ}.) Before (a time)? (\emph{zhīqián}.) After (a time)? (\emph{zhīhòu}.) Morning and evening? (\textbf{\zh{早晚}}, \emph{zǎowǎn}.) More often than not? (\emph{wǎngwǎng}.) Scissors? (\emph{jiǎndāo}.) A rubbish bin? (\emph{lājītǒng}.) A plastic bag? (\emph{sùliàodài}.) A bottle? (\emph{píngzi}.) A drawer? (\emph{chōuti}.) To smash? (\emph{dǎpò}.) To have a fall? (\emph{shuāi}.) To regret? (\emph{hòuhuǐ}.) At that time? (\emph{dāngshí}.) In the end? (\emph{jiéguǒ}.)
\end{tcolorbox}
\section[(dialogue)]{What a Day — once more}

\label{lesson:ZH-R187-review-final-187}



\begin{quote}
Say the Mandarin for to smash and to have a fall.
\end{quote}

\subsection*{Your turn}
\begin{itemize}
\raggedright
  \item \emph{Say it:} to smash, to have a fall, to regret, at that time, in the end
\end{itemize}

\begin{tcolorbox}[breakable,colback=teal!4,colframe=teal!35!black,title={Before you move on}]
To smash? (\emph{dǎpò}.) To have a fall? (\emph{shuāi}.) To regret? (\emph{hòuhuǐ}.) At that time? (\emph{dāngshí}.) In the end? (\emph{jiéguǒ}.)
\end{tcolorbox}
